%---------- Inleiding ---------------------------------------------------------

% TODO: Is dit voorstel gebaseerd op een paper van Research Methods die je
% vorig jaar hebt ingediend? Heb je daarbij eventueel samengewerkt met een
% andere student?
% Zo ja, haal dan de tekst hieronder uit commentaar en pas aan.

%\paragraph{Opmerking}

% Dit voorstel is gebaseerd op het onderzoeksvoorstel dat werd geschreven in het
% kader van het vak Research Methods dat ik (vorig/dit) academiejaar heb
% uitgewerkt (met medesturent VOORNAAM NAAM als mede-auteur).
% 

\section{Inleiding}%
\label{sec:inleiding}

Goossens Warmte- en Koeltechniek is een Belgiesch HVAC bedrijf dat Installware gebruikt ter ondersteuning van de buitendienst.
Deze bachelorproef richt zich op de registratie en verwerking van digitale werkbongegevens binnen deze bedrijfscontext.
De doelgroep voor deze bachelorproef zijn de buitendienstmedewerkers en de kantoormedewerkers die deze gegevens
raadplegen en verwerken.

Wanneer medewerkers gegevens aanpassen zonder betrouwbare netwerkverbinding,
moet duidelijk zijn hoe deze gegevens lokaal worden bewaard en hoe ze worden
verwerkt wanneer ze terug netwerkverbinding hebben.

Het onderzoek vergelijkt twee concreet uitgewerkte aanpakken:
Een lokale wachtrij met centrale conflictcontrole en een aanpak met conflictvrije gerepliceerde datatypes.
Het wordt afgebakend tot werknotities en materiaalregistraties. Tijdens de analyse van de huidige werkwijze
wordt getoetst of deze gegevens en de bijbehorende gebruiksscenario’s representatief zijn voor het bedrijf.
De technische vergelijking die gemaakt zal worden bij de proof of concept zal geen toegang tot de
broncode van Installware nodig hebben.

De hoofdonderzoeksvraag: “Welke van twee synchronisatieaanpakken, een lokale wachtrij met centrale conflictcontrole of conflictvrije gerepliceerde datatypes gebaseerde samenvoeging, biedt de beste afweging tussen gegevenscorrectheid, synchronisatietijd en implementatiecomplexiteit voor offline bewerkte werknotities en materiaalregistraties op digitale werkbonnen bij Goossens Warmte- en Koeltechniek?”

Deelvragen probleemdomein:
(1) Hoe registreren en verwerken buitendienstmedewerkers en kantoormedewerkers momenteel werknotities en materiaalregistraties met Installware?
(2) Welke mogelijkheden en beperkingen ervaren zij bij het raadplegen, bewerken en synchroniseren van deze gegevens tijdens en na een onderbroken netwerkverbinding?
(3) In welke situaties wijzigen meerdere medewerkers dezelfde werkbon en bestaan daar al bepaalde regels voor?
(4) Welke functionele en technische vereisten stelt het bedrijf aan offline gegevensvastlegging en de daaropvolgende synchronisatie?

Deelvragen oplossingsdomein:

(1) Hoe kunnen beide synchronisatieaanpakken worden ingericht voor de geselecteerde werkbongegevens en bedrijfsvereisten?
(2) Hoe verwerken beide implementaties wijzigingen bij verbindingsuitval, onderbroken synchronisatie, herhaalde verzending en gelijktijdige bewerking?
(3) Hoe verschillen beide implementaties in gegevenscorrectheid en synchronisatietijd bij dezelfde gegevens, gebruiksscenario’s en netwerkomstandigheden?
(4) Welke implementatiecomplexiteit en beperkingen brengt elke aanpak met zich mee?

Het doel van deze bachelorproef is een onderbouwd advies te formuleren over de geschiktheid van beide synchronisatieaanpakken
voor de bedrijfsvereisten. Het onderzoek levert een analyse van de huidige werkwijze en vereisten, een zelfstandig prototype met
twee vergelijkbare implementaties en een evaluatie op basis van reproduceerbare tests op.
Gegevenscorrectheid wordt daarbij beoordeeld aan de hand van zowel het behoud van registraties als de naleving
van de vastgestelde regels voor werkbonnen. Het onderzoek heeft niet als doel een volledige vervanging van Installware,
maar een technisch onderbouwde vergelijking die het bedrijf inzicht geeft in de mogelijkheden en beperkingen van beide aanpakken.

%---------- Stand van zaken ---------------------------------------------------

\section{Literatuurstudie}%
\label{sec:literatuurstudie}

Bij mobiele applicaties voor buitendienstmedewerkers is het belangrijk dat de ingevoerde gegevens betrouwbaar zijn.
\textcite{KnowlesEtAl2014} onderzoeken dit bij medewerkers in het wegenonderhoud. \textcite{HardingEtAl2019} onderzoeken
hoe medewerkers in transportonderhoud gegevens gebruiken bij hun werk. Daarbij brengen ze de bestaande werkprocessen in kaart
via interviews en observaties. Deze onderzoeken gaan niet specifiek over HVAC-bedrijven, maar tonen wel waarom het belangrijk is om
eerst te begrijpen hoe medewerkers gegevens registreren en gebruiken. In deze bachelorproef wordt daarom eerst de huidige
werkwijze van het bedrijf onderzocht.

Om ook zonder netwerkverbinding gegevens te kunnen raadplegen, kan een mobiele applicatie gegevens lokaal opslaan.
De Android-documentatie beschrijft hoe een offline-first applicatie hiervoor een lokale gegevensbron combineert met
een gegevensbron op het netwerk. Wanneer gebruikers ook offline gegevens moeten kunnen aanpassen, kunnen wijzigingen
eerst lokaal worden bewaard en daarna in een wachtrij worden geplaatst. Zodra er opnieuw verbinding is, probeert de
applicatie de wijzigingen door te sturen~\autocite{Google2026OfflineFirst}. Ook \textcite{NingrumEtAl2026} bespreken
lokale opslag en uitgestelde synchronisatie als veelgebruikte onderdelen van applicaties met offlineondersteuning.

Er bestaan verschillende manieren om voor synchronisatie in mobiele applicaties te zorgen.
\textcite{McCormickSchmidt2012} beschrijven onder meer wanneer gegevens worden gesynchroniseerd,
welke gegevens lokaal beschikbaar blijven en welke wijzigingen worden doorgestuurd. De geschikte
aanpak hangt vooral af van de netwerkverbinding en van hoe recent de gegevens moeten zijn.
Een wachtrij kan wijzigingen bewaren totdat ze kunnen worden verstuurd, maar lost niet automatisch conflicten op.
Een conflict kan bijvoorbeeld ontstaan wanneer twee medewerkers dezelfde gegevens aanpassen terwijl ze offline zijn.
Er zijn dan extra regels nodig om te bepalen welke wijzigingen worden behouden. In de Android-documentatie worden
onder andere het bijhouden van versies en last-write-wins besproken. Bij last-write-wins wordt één wijziging als de laatste
beschouwd en behouden, waardoor een andere wijziging kan worden overschreven~\autocite{Google2026OfflineFirst}.

Een andere aanpak maakt gebruik van conflictvrije gerepliceerde datatypes, ook CRDT's genoemd. Daarmee kunnen verschillende
apparaten elk hun lokale kopie van gegevens aanpassen. Wanneer alle kopieën dezelfde wijzigingen hebben ontvangen, komen ze 
onder de voorwaarden van het gekozen CRDT tot dezelfde toestand~\autocite{PreguicaEtAl2018}. Er bestaan verschillende CRDT's,
bijvoorbeeld voor tellers, verzamelingen en tekstgegevens. Elk type heeft regels die bepalen hoe wijzigingen worden samengevoegd.
De keuze van het datatype en de regels blijft dus belangrijk. Het automatisch samenvoegen van wijzigingen betekent niet dat
het resultaat ook altijd overeenkomt met wat de gebruikers bedoelden~\autocite{PreguicaEtAl2018}.

\textcite{KleppmannEtAl2019} bespreken hoe lokale gegevensopslag en samenwerking tussen apparaten gecombineerd kunnen worden.
Ze noemen dit local-first software en onderzoeken hoe CRDT's hierbij kunnen helpen. Local-first gaat verder dan
alleen offline werken en houdt bijvoorbeeld ook rekening met de controle van gebruikers over hun gegevens.
Een wachtrij en CRDT's zijn daarbij geen volledig
aparte technieken. Ook een CRDT-gebaseerde applicatie moet wijzigingen bewaren en doorsturen. De vergelijking in dit onderzoek
gaat daarom over twee concrete implementaties die wijzigingen en conflicten op een andere manier
verwerken~\autocite{PreguicaEtAl2018}.

De literatuur beschrijft verschillende mogelijkheden, maar geeft nog geen antwoord op welke aanpak het meest geschikt is
voor de werkbonnen van het onderzochte bedrijf. Dat apparaten na synchronisatie dezelfde gegevens bevatten, is bovendien niet
voldoende. Het resultaat moet ook voldoen aan de afspraken over werknotities en materiaalregistraties. Daarom worden de
bedrijfsvereisten gebruikt om te bepalen wat als een correct resultaat wordt beschouwd. Beide implementaties worden vervolgens
getest met dezelfde gegevens en situaties. Zo wordt onderzocht hoe ze verschillen in gegevenscorrectheid, synchronisatietijd
en implementatiecomplexiteit.

%---------- Methodologie ------------------------------------------------------
\section{Methodologie}%
\label{sec:methodologie}

Dit onderzoek bestaat uit een analyse van de bedrijfscasus en een vergelijkende studie met een proof of concept.
Eerst worden de huidige werkwijze en de bedrijfsvereisten onderzocht. Daarna worden twee implementaties
uitgewerkt en onder dezelfde omstandigheden getest. De resultaten vormen de basis voor een advies aan het bedrijf.

In de eerste fase wordt gekeken hoe werknotities en materiaalregistraties momenteel met Installware worden verwerkt.
Om dit achter te halen wordt contact opgenomen met co-promotor Jarne Schagen. De vragen gaan over de gegevens die medewerkers
registreren, de mogelijkheden tijdens netwerkuitval, de stappen
na verbindingsherstel en situaties waarin meerdere medewerkers dezelfde werkbon aanpassen.

Het proof of concept bestaat uit een eenvoudige mobiele testapplicatie met lokale opslag en een backend.
Als voorlopig technisch uitgangspunt wordt React Native met Expo en
TypeScript gebruikt~\autocite{ReactNativeTypeScript}, met SQLite voor lokale opslag~\autocite{ExpoSQLite}. Voor de
CRDT-gebaseerde implementatie wordt Yjs gebruikt~\autocite{YjsDocumentation}. De eerste implementatie bewaart wijzigingen in een lokale
wachtrij. De backend controleert via versienummers of een wijziging op verouderde gegevens is gebaseerd. Conflicten worden
zichtbaar gemaakt en volgens de vooraf vastgelegde regels afgehandeld. De tweede implementatie gebruikt CRDT's om wijzigingen
samen te voegen. Beide varianten ondersteunen dezelfde handelingen voor werknotities en materiaalregistraties. Ze gebruiken
zoveel mogelijk dezelfde gebruikersinterface, gegevens en testomgeving, zodat er zo weinig mogelijk verschillen zijn die invloed zouden kunnen hebben.

De tests gebruiken fictieve werkbonnen die aansluiten bij de vastgestelde bedrijfsvereisten. Er worden scenario's opgesteld voor
herstel van de verbinding, een nieuwe onderbreking tijdens synchronisatie, herhaalde verzending en gelijktijdige
wijzigingen door meerdere gebruikers.

Gegevenscorrectheid wordt beoordeeld door de gegevens te vergelijken met de resultaten die vooraf werden vastgelegd.
Daarbij worden ontbrekende registraties, ongewenste dubbels, afwijkingen van de werkbonregels en verschillen tussen gegevenskopieën
bijgehouden. De synchronisatietijd wordt gemeten
vanaf het herstel van de verbinding tot alle testwijzigingen zijn verwerkt. Elk testgeval wordt dertig keer uitgevoerd vanuit dezelfde
beginsituatie. De resultaten worden samengevat met de mediaan en de spreiding.

Tot slot worden beide implementaties naast de bedrijfsvereisten gelegd. Het advies vermeldt zowel de voordelen als de beperkingen
van elke aanpak.

%---------- Verwachte resultaten ----------------------------------------------
\section{Verwacht resultaat, conclusie}%
\label{sec:verwachte_resultaten}

Het verwachte resultaat is een vergelijking van twee synchronisatieaanpakken
voor werknotities en materiaalregistraties, ondersteund door een proof of
concept. De vergelijking geeft inzicht in gegevenscorrectheid,
synchronisatietijd en implementatiecomplexiteit.

Er wordt verwacht dat de verschillen vooral zichtbaar gaan zijn bij
gelijktijdige wijzigingen. De wachtrijaanpak kan hierbij expliciete
conflictcontrole vragen, terwijl CRDT's bepaalde wijzigingen automatisch
kunnen samenvoegen.

De verwachte conclusie is een onderbouwd advies over welke aanpak het
beste aansluit bij de bedrijfsvereisten en welke beperkingen er zijn.
Dit geeft het bedrijf inzicht voor toekomstige keuzes rond
digitale werkbonnen.